\documentclass[journal, a4paper]{IEEEtran}

\usepackage[T1]{fontenc}
\usepackage[utf8]{inputenc}
\usepackage[spanish, es-noshorthands]{babel}
\usepackage{amsmath, amssymb, bm}
\usepackage{graphicx}
\usepackage{booktabs}
\usepackage{tabularx}
\usepackage{ragged2e}
\usepackage[american]{circuitikz}
\usepackage{pgfplots}
\pgfplotsset{compat=1.18}
\usepackage{hyperref}

\renewcommand{\tablename}{Tabla}

\begin{document}

\title{Verificación Experimental de Superposición y Equivalentes de Thévenin/Norton en una Red Resistiva}

\author{
    [Nombre del Estudiante 1], [Nombre del Estudiante 2]\\
    Ingeniería Mecatrónica, Universidad Católica Boliviana ``San Pablo'' -- Sede Tarija\\
    IMT-121 --- Circuitos Electrónicos I\\
    Fecha: \today \quad | \quad Grupo: [ID del Grupo]
}

\maketitle

\begin{abstract}
% Escriba un párrafo (80-120 palabras) indicando el objetivo, el método, la comparación principal y la conclusión técnica.
[Escribir resumen aquí]
\end{abstract}

\begin{IEEEkeywords}
Superposición, equivalente de Thévenin, equivalente de Norton, linealidad, LTSpice, validación experimental.
\end{IEEEkeywords}

\section{Introducción y Objetivo}
% Explique la relación que se evalúa (linealidad, equivalencia de puertos) y el objetivo explícito del experimento.
[Escribir introducción aquí]

\section{Modelo del Circuito y Predicciones Analíticas}

\subsection{Definición del Circuito y Referencias}
% Incluya el esquemático usando circuitikz del circuito base indicando las referencias de tensión V_a, el puerto a-b, las polaridades y los valores nominales.
[Insertar esquemático aquí]

\subsection{Relaciones Fundamentales}
Las respuestas parciales bajo el principio de superposición se combinan algebraicamente:
\begin{equation}
    V_{a,\text{total}} = V_{a}^{(V1)} + V_{a}^{(V2)}
\end{equation}
Los parámetros del equivalente se relacionan mediante:
\begin{align}
    V_{th} &= V_{oc} \\
    R_{th} &= R_{\text{vista desde puerto}} \\
    I_{N} &= \frac{V_{th}}{R_{th}}
\end{align}

\subsection{Predicciones Analíticas}
\begin{table}[htbp]
    \centering
    \caption{Predicciones Analíticas Base}
    \label{tab:predicciones_a}
    \begin{tabularx}{\columnwidth}{>{\RaggedRight\arraybackslash}X >{\centering\arraybackslash}X c}
        \toprule
        \textbf{Magnitud} & \textbf{Predicción Analítica} & \textbf{Unidad} \\
        \midrule
        $V_{a}^{(V1)}$ & & V \\
        $V_{a}^{(V2)}$ & & V \\
        $V_{oc} = V_{th}$ & & V \\
        $R_{th}$ & & k$\Omega$ \\
        $I_{N}$ & & mA \\
        $V_L$ ($R_L$ conectada) & & V \\
        \bottomrule
    \end{tabularx}
\end{table}

\section{Metodología}
\subsection{Simulación}
% Describa la configuración en LTSpice para obtener las contribuciones de superposición, el voltaje a circuito abierto y el comportamiento con carga.
[Descripción de simulación]

\subsection{Montaje Experimental}
% Documente la implementación en protoboard, configuración de fuentes, instrumentos y nodos de referencia.
[Descripción del montaje]

\subsection{Procedimiento}
% Secuencia numerada de medición de contribuciones, V_oc, R_th y validación con carga.
% ATENCIÓN: Indique explícitamente que I_N se obtuvo analíticamente o por simulación, NO mediante un corto físico del puerto, a menos que el instructor lo haya autorizado expresamente.
[Secuencia de pasos]

\section{Resultados}

\begin{table*}[htbp]
    \centering
    \caption{Comparación de Evidencia Experimental y Simulada}
    \label{tab:resultados_a}
    \renewcommand{\arraystretch}{1.3}
    \begin{tabularx}{\textwidth}{>{\RaggedRight\arraybackslash}X *{4}{>{\centering\arraybackslash}X}}
        \toprule
        \textbf{Magnitud} & \textbf{Teoría} & \textbf{SPICE} & \textbf{Medición} & \textbf{Error Relativo (\%)} \\
        \midrule
        $V_a^{(V1)}$ & & & & \\
        $V_a^{(V2)}$ & & & & \\
        $V_{a,\text{total}}$ & & & & \\
        $V_{oc}$ & & & & \\
        $R_{th}$ & & & & \\
        $I_{N}$ & & & \textit{N/A (Cálculo)} & \textit{N/A} \\
        $V_L$ (con $R_L$) & & & & \\
        $I_L$ (con $R_L$) & & & & \\
        \bottomrule
    \end{tabularx}
\end{table*}

\section{Discusión}
% Responda explícitamente:
% 1. ¿Se sostuvo experimentalmente la superposición?
% 2. ¿El equivalente de Thévenin reprodujo el voltaje terminal de la red original bajo carga?
% 3. Justifique físicamente la causa más probable de las discrepancias observadas (observación -> hipótesis -> evidencia).
[Escribir discusión aquí]

\section{Conclusiones}
% Proporcione de 3 a 5 conclusiones técnicas precisas y ligadas a su evidencia cuantitativa.
[Escribir conclusiones aquí]

\begin{thebibliography}{00}
\bibitem{guia_a} Instructor del Curso, \emph{Experiencia Integrada A — Guía de Laboratorio}, Universidad Católica Boliviana, 2026.
\end{thebibliography}

\end{document}
